\documentclass[10.5pt,a4paper]{article}

\usepackage[utf8]{inputenc}
\usepackage{CJKutf8}
\usepackage[
  top=0.80cm,
  bottom=0.90cm,
  left=1.25cm,
  right=1.25cm
]{geometry}
\usepackage{enumitem}
\usepackage{titlesec}
\usepackage{hyperref}
\usepackage{array}
\usepackage{tabularx}
\usepackage{graphicx}
\usepackage{fontawesome5}

\pagestyle{empty}
\setlength{\parindent}{0pt}
\setlength{\parskip}{0pt}
\linespread{1.02}

\hypersetup{
  colorlinks=true,
  linkcolor=black,
  urlcolor=blue
}

\setlist[itemize]{
  leftmargin=1.25em,
  itemsep=0.12em,
  topsep=0.18em,
  parsep=0em,
  partopsep=0em
}

\renewcommand\labelitemi{\scriptsize$\bullet$}

\titleformat{\section}
  {\large\bfseries}
  {}
  {0em}
  {}
  [\vspace{-0.08em}\titlerule]

\titlespacing*{\section}{0pt}{0.75em}{0.45em}

% =========================
% 自定义命令
% =========================
\newcommand{\resumeItem}[1]{\item #1}

\newcommand{\projectHeader}[3]{%
  \noindent
  \begin{tabular*}{\textwidth}{@{\extracolsep{\fill}}lr}
    \textbf{#1} &
    \textbf{#2}\hspace{1.5em}#3
  \end{tabular*}

  \vspace{0.15em}
}

\newcommand{\eduEntry}[6]{%
\textbf{#1}
\hspace{1.0em}
\textbf{#2}
\hspace{1.0em}
\textbf{#3}
\hfill
#4

\vspace{0.25em}

#5

\vspace{0.15em}

#6

\vspace{0.5em}
}

\newcommand{\subTitle}[2]{
  \textbf{#1}
  \hfill
  #2
}

\begin{document}
\begin{CJK*}{UTF8}{gbsn}

\newcommand{\contactItem}[2]{
  \makebox[1.2em][c]{#1}\hspace{0.25em}#2
}

% =========================
% 顶部个人信息
% =========================
\noindent
\begin{tabularx}{\textwidth}{@{}X@{\hspace{0.8cm}}p{2.35cm}@{}}
  \begin{minipage}[c][2.85cm][c]{\linewidth}
    \centering

    {\hspace*{-1.0cm}\huge \textbf{陈春媚}} \\[0.6em]
%    {\hspace*{-1.0cm}\Large \textbf{求职意向：软件工程师、软件测试工程师}} \\[0.55em]

    \footnotesize
    \begin{tabular}{@{}r@{\hspace{0.35em}}l@{\hspace{7.5em}}r@{\hspace{0.35em}}l@{}}
      \Large \faPhone* & \Large
      15878562962
      &
      \Large \faEnvelope & \Large
      \href{mailto:1476137517@qq.com}{1476137517@qq.com}
      \\[0.25em]

      \Large \faGithub & \Large
      \href{https://github.com/chenchunmei-daphne}{chenchunmei-daphne}
      &
      \Large \faHome & \Large
      现居湖南湘潭｜籍贯广西

    \end{tabular}
  \end{minipage}
  &
  \begin{minipage}[c][2.85cm][c]{2.35cm}
    \raggedleft
    \includegraphics[width=2.25cm,height=2.85cm]{photo.jpg}
  \end{minipage}
\end{tabularx}

\vspace{0.15em}
% =========================
% 教育经历
% =========================
\section*{教育经历}

\eduEntry
{湘潭大学}
{硕士     }
{应用统计}
{2024.09--2027.06}
{研究方向：\textbf{深度学习}，\textbf{信道矩阵重构}，PINN 求解偏微分方程}
{相关课程：机器学习，最优化理论与算法，多元统计分析，时间序列分析}

\eduEntry
{福州大学}
{本科    }
{数学与应用数学}
{2020.09--2024.06}
{相关课程：数学分析，高等代数，数据结构与算法，数据库系统(SQL)，数学软件(MATLAB)}


% =========================
% 项目经历
% =========================
\section*{项目经历}

\projectHeader
{华为合作项目｜基于算子学习的信道重构}
{主要开发人员}
{2025.07--2026.09}

\smallskip

\textbf{项目概述：}{参与与华为合作的算子学习信道重构项目，面向大规模\ MIMO 不完全观测信道重构问题，负责数据生成、模型开发、训练测试与实验迭代。}

\par \medskip
\textbf{主要工作与结果：}
\begin{itemize}

\resumeItem{\textbf{基于\ QuaDRiGa 搭建无线信道数据生成与预处理流程}，完成\ LOS、NLOS 等场景的数据生成和预处理工作。}

\resumeItem{\textbf{基于\ PyTorch 设计并实现\ Transformer 信道重构模型及训练框架}，依托超算互联网平台完成约\ \textbf{8 万样本的多\ GPU 并行训练}，测试集\textbf{频谱效率较基线提升约\ 30\%}。}

\resumeItem{\textbf{负责模型与实验流程迭代}，围绕数据归一化、损失函数和训练稳定性开展实验，完善检查点与指标保存机制，提升训练稳定性及实验可复现性。}

\end{itemize}

\vspace{0.7em}

\projectHeader
{FEALPy 科学计算\ RAG 智能体}
{核心开发人员}
{2026.08--至今}

\smallskip

\textbf{项目概述：}{面向\ FEALPy 新成员开发领域智能助手，将自然语言计算需求转换为可调用的\ \texttt{bm.*} 接口及示例代码，构建涵盖意图理解、知识检索、答案生成、事实校验与多轮会话的智能体工作流。}

\par \medskip
\textbf{主要工作与成果：}
\begin{itemize}

\resumeItem{\textbf{设计并实现领域智能体工作流}，使用通义千问完成结构化意图解析与答案生成，串联意图识别、知识检索、候选重排、答案校验和会话记忆模块，支持后端、数据类型、设备及张量形状等约束的自然语言理解。}

\resumeItem{\textbf{构建\ FEALPy 接口\ RAG 知识库}，整理\ 220 条多后端接口数据，使用\ BGE-M3 生成\ 1024 维向量并通过\ Chroma 持久化，结合语义召回、关键词匹配与分类软加权完成候选接口排序。}

\resumeItem{\textbf{实现基于检索证据的幻觉防护与降级机制}，校验回答中的\ FEALPy、NumPy 与\ PyTorch API 是否来源于检索结果；当模型调用、结构化输出或事实校验异常时，自动回退至规则解析与模板回答。}

\resumeItem{\textbf{完成多轮会话、服务化与评测闭环}，基于\ session\_id 隔离上下文，使用\ FastAPI 与\ Gradio 提供聊天、检索及接口查询功能；构建\ 55 条领域问题评测集，通过\ Recall@K、答案准确率及消融实验定位检索瓶颈。}

\end{itemize}

\vspace{0.7em}

\projectHeader
{FEALPy 深度学习模块与示例工程开发}
{主要开发人员}
{2025.06--2025.08}

\smallskip

\textbf{项目概述：}{参与\href{https://github.com/chenchunmei-daphne/fealpy}{\ FEALPy} 开源数值计算软件开发，使用\ Python、PyTorch 开发\ PINN、PENN、RFM 深度学习模块及偏微分方程求解示例。}

\par \medskip
\textbf{主要工作与结果：}
\begin{itemize}

\resumeItem{\textbf{设计并实现\ PINN、PENN 与\ RFM 模块化代码框架}，封装模型构建、训练、损失计算、优化器配置及结果可视化等功能。}

\resumeItem{\textbf{开发多类\ PDE 求解示例程序}，完成模型训练、数值实验、结果验证及程序调试，提高模块的复用性和可维护性。}

\end{itemize}

\vspace{0.7em}

% =========================
% 技术影响力
% =========================
\section*{技术影响力}

\subTitle{开源贡献}{2024.09--至今}

\begin{itemize}

\resumeItem{参与
\href{https://github.com/chenchunmei-daphne/fealpy}{\ FEALPy}
开源软件开发，主要贡献\textbf{深度学习\ PDE 求解、有限差分模块及相关示例}；熟悉\ Git/GitHub 协作，累计提交\ \textbf{PR 160+ 次}，贡献代码\ \textbf{6000+ 行}。}

\end{itemize}

% =========================
% 专业能力
% =========================
\section*{专业能力}

\begin{itemize}

\resumeItem{\textbf{编程开发：}熟悉\ Python、PyTorch、Linux 与\ Git，具备模块开发、数据处理、程序调试及问题定位能力。}

\resumeItem{\textbf{深度学习：}熟悉\ Transformer、CNN、PINN 等模型，能够完成模型构建、训练与性能分析，具备多\ GPU 并行及边读取边训练实践经验。}

\resumeItem{\textbf{大模型与智能体应用：}具备\ RAG 领域智能体开发经验，熟悉提示词设计、结构化输出、向量检索、多轮记忆、幻觉校验与失败回退，使用过\ BGE-M3、Chroma、通义千问、FastAPI 与\ Gradio。}

\resumeItem{\textbf{技能证书：}全国计算机等级考试（NCRE）二级，大学英语四级。}


\end{itemize}

\end{CJK*}
\end{document}
